% Onda trasversale (corda) e onda longitudinale (suono): le particelle del
% mezzo oscillano perpendicolari o parallele alla direzione di propagazione.
\begin{tikzpicture}[>=Stealth, font=\small]
  % (a) onda trasversale
  \node[anchor=west] at (-0.6,1.15) {(a) onda trasversale};
  \draw[formula, very thick, domain=0:8, samples=120] plot (\x, {0.55*sin(\x*90)});
  \foreach \x in {1,3,5,7} {
    \draw[<->, esempio, thick] (\x,-0.75) -- (\x,0.75);
  }
  \draw[->, thick] (8.4,0) -- (9.6,0) node[right, align=left] {propagazione};
  \node[esempio, align=center] at (2,-1.0) {oscillazione delle particelle};
  % (b) onda longitudinale
  \begin{scope}[yshift=-3.3cm]
    \node[anchor=west] at (-0.6,1.15) {(b) onda longitudinale};
    \foreach \i in {0,...,50} {
      \pgfmathsetmacro{\xo}{\i*0.16}
      \pgfmathsetmacro{\xx}{\xo + 0.26*sin(\xo*90)}
      \draw[formula, thick] (\xx,-0.5) -- (\xx,0.5);
    }
    \foreach \x in {2.0,6.0} {
      \draw[<->, esempio, thick] ({\x-0.35},0.75) -- ({\x+0.35},0.75);
    }
    \draw[->, thick] (8.4,0) -- (9.6,0) node[right] {propagazione};
    \node[below] at (2,-0.55) {compressione};
    \node[below] at (6,-0.55) {compressione};
    \node[below] at (4,-0.55) {rarefazione};
    \draw[|<->|] (2,-1.3) -- (6,-1.3) node[midway, below] {$\lambda$};
  \end{scope}
\end{tikzpicture}
